\clearpage\thispagestyle{empty}
\pdfbookmark[0]{Source for the 2 April 1808 Notice on the Comet of 1807}{comet1807-apr1808-sources}
\begin{center}\textbf{Source for the 2 April 1808 Notice on the Comet of 1807}\end{center}
\noindent This AI translation follows the notice dated 2 April 1808 in \href{https://gdz.sub.uni-goettingen.de/id/PPN236007467}{Gauss’s \textit{Werke}, vol.~VI} (Göttingen, 1874), pp.~300--301. The notice appeared under the heading \textit{Göttingen} in \textit{Göttingische gelehrte Anzeigen}, issue~53, \href{https://manifests.sub.uni-goettingen.de/iiif/image/gdz:PPN319721507_1808:00000525/full/1600,/0/default.jpg}{pp.~521}, \href{https://manifests.sub.uni-goettingen.de/iiif/image/gdz:PPN319721507_1808:00000526/full/1600,/0/default.jpg}{522} and \href{https://manifests.sub.uni-goettingen.de/iiif/image/gdz:PPN319721507_1808:00000527/full/1600,/0/default.jpg}{523}. The notice presents Bessel’s observations from Lilienthal and a separate observation by Olbers from Bremen. Its reporter refers to Gauss in the third person.
\par\medskip
The collected edition uses decimal points in Bessel’s observations and decimal commas in the tables of parabolic and elliptical elements. This distinction is retained here. For related English terminology, see Charles Henry Davis’s 1857 translation of Gauss’s \textit{Theory of the Motion of the Heavenly Bodies Moving About the Sun in Conic Sections}, \href{https://www2.karlin.mff.cuni.cz/~tuma/Aplikace21/theoryofmotionof00gausuoft.pdf#page=89}{Book~I, Section~II, Article~57, p.~65}. Translation and typesetting produced by AI.
\clearpage
